%\chapter*{附錄：教師使用指南}
%\addcontentsline{toc}{chapter}{附錄：教師使用指南}

%\appendix
\renewcommand{\appendixname}{附錄}
\renewcommand{\thechapter}{\arabic{chapter}}
\renewcommand{\thefigure}{\appendixname~\thechapter.\arabic{figure}}  % Appendix 1.1, 1.2, …

\refstepcounter{chapter}
\chapter*{附錄：教師使用指南}
\addcontentsline{toc}{chapter}%
  {\appendixname~\thechapter\quad 教師使用指南}



此書適合成人主日學，或是小組研讀。以下建議進度表，方便教師或是帶領人安排進度。我的建議是採用八週作為主日學正式課程，並保留十週作為深入研讀版本。 

\section*{一、六週：掌握啟示錄的主線}

適合時間有限、希望在六週內掌握全書主要信息的班級。每週閱讀一個或數個相關章節，重點放在基督的身分、救贖、審判、得勝與新創造。
{
\renewcommand{\arraystretch}{1}
\begin{longtable}{
  >{\centering\arraybackslash}p{0.5cm}
  p{2cm}
  p{1cm}
  p{1.5cm}
  p{4cm}
}
週 &
主題 &
本書 &
啟示錄 &
教學重點 \\
\hline
\endfirsthead

週 &
主題 &
本書 &
啟示錄 &
教學重點 \\
\hline
\endhead

一 &
認識榮耀的基督與祂的教會 &
1-2 &
1-3 &
認識復活、掌權的基督；了解七教會的稱許、責備、命令與應許，思想教會如何忠心持守信仰。 \\

二 &
榮耀的寶座與羔羊展開書卷 &
3-4 &
4-5 &
從創造者的敬拜到被殺羔羊的敬拜，認識基督配得展開書卷，並掌管歷史。 \\

三 &
羔羊掌管歷史與審判 &
4-5 &
6-7 &
認識七印、四騎士、殉道者的呼求，以及受印的僕人和萬國萬民的敬拜；在苦難中持守盼望。 \\

四 &
基督審判世界，祂的百姓忠心作見證 &
5-6 &
8-11 &
認識七號的審判、兩個見證人和第七號，思想神的審判、憐憫與教會的見證使命。 \\

五 &
基督勝過龍、獸與巴比倫 &
7-10 &
12-19 &
認識龍、獸、假先知和巴比倫所代表的敵對勢力；思想聖徒如何靠羔羊得勝，並等候基督再來。 \\

六 &
基督擊敗撒但，神帶領百姓進入榮耀 &
11-12 &
20-22 &
認識千禧年不同觀點、最後審判、新天新地與新耶路撒冷；以「主耶穌啊，我願你來」作為課程的盼望與回應。 \\

\end{longtable}
}

\section*{二、八週：兼顧全書概覽與重要經文}

這是我最推薦的方案。每週可以安排約 60-75 分鐘，既能介紹啟示錄的主要異象，也能討論教會生活、苦難中的見證、末世觀點及基督徒的終極盼望。
{
\renewcommand{\arraystretch}{1}
\begin{longtable}{
  >{\centering\arraybackslash}p{0.5cm}
  p{2cm}
  p{1cm}
  p{1.5cm}
  p{4cm}
}
週 &
主題 &
本書 &
啟示錄 &
教學重點 \\
\hline
\endfirsthead

週 &
主題 &
本書 &
啟示錄 &
教學重點 \\
\hline
\endhead

一 &
認識榮耀的基督 &
1 &
1:1-20 &
了解約翰與拔摩島的背景；從基督的受苦、復活與掌權，認識啟示錄的中心人物。 \\

二 &
基督對祂的教會說話 &
2 &
2:1-3:22 &
研讀七教會的稱許、責備、命令與應許；思想教會如何悔改、持守真道，並忠心到底。 \\

三 &
榮耀的寶座與羔羊 &
3 &
4:1-5:14 &
從創造萬有的主到被殺的羔羊，認識天上敬拜的中心，以及基督為何配展開書卷。 \\

四 &
羔羊展開書卷 &
4 &
6:1-7:17 &
研讀七印、四騎士、殉道者、受印的僕人與萬國萬民的敬拜；思想苦難中的忍耐與救恩的盼望。 \\

五 &
基督審判世界，教會忠心作見證 &
5-6 &
8:1-11:19 &
研讀七號、兩個見證人和第七號；理解審判中的限制、人的不悔改，以及教會的見證使命。 \\

六 &
基督勝過龍與獸的權勢 &
7 &
12:1-14:20 &
了解婦人、男孩子、龍、海獸與假先知的象徵；思想聖徒如何靠羔羊的血和所見證的道得勝。 \\

七 &
公義審判與羔羊的得勝 &
8-10 &
15:1-19:21 &
研讀七碗、巴比倫傾倒、羔羊婚筵與騎白馬的基督；比較世界的權勢與基督最終的得勝。 \\

八 &
基督擊敗撒但，被恩典帶入榮耀 &
11-12 &
20:1-22:21 &
研讀千禧年、最後審判、新天新地與新耶路撒冷；認識不同末世觀點，並以盼望等候主再來。 \\

\end{longtable}
}


\section*{三、十週：深入研讀}

十週方案最能發揮本書的特色，尤其適合主日學教師培訓、成人查經班，以及希望進一步理解啟示錄象徵與末世觀點的讀者。

{
\renewcommand{\arraystretch}{1}
\begin{longtable}{
  >{\centering\arraybackslash}p{0.5cm}
  p{2cm}
  p{1cm}
  p{1.5cm}
  p{4cm}
}
週 &
主題 &
本書 &
啟示錄 &
教學重點 \\
\hline
\endfirsthead

週 &
主題 &
本書 &
啟示錄 &
教學重點 \\
\hline
\endhead

一 &
認識榮耀的基督 &
1 &
1:1-20 &
了解啟示錄的背景、約翰的見證，以及基督從受苦到得榮耀的身分；建立以基督為中心的閱讀方向。 \\

二 &
基督對祂的教會說話 &
2 &
2:1-3:22 &
逐一認識七教會；比較各教會的處境、優點、問題與應許，反思今日教會的信仰生活。 \\

三 &
榮耀的寶座與羔羊 &
3 &
4:1-5:14 &
研讀天上的寶座、四活物、二十四位長老與羔羊；理解創造與救贖如何匯合於敬拜。 \\

四 &
羔羊展開書卷 &
4 &
6:1-7:17 &
深入研讀七印、四騎士、殉道者的呼求、十四萬四千受印者與萬國萬民；比較七印的連續發生與重述觀點。 \\

五 &
七號與神的審判 &
5 &
8:1-9:21 &
研讀前六號及相關異象；注意「三分之一」所呈現的限制，並探討人為何在災難中仍不悔改。 \\

六 &
在苦難中忠心作見證 &
6 &
10:1-11:19 &
研讀小書卷、兩個見證人與第七號；理解神的話語、教會的見證，以及基督國度的宣告。 \\

七 &
基督勝過龍與獸的權勢 &
7 &
12:1-14:20 &
從婦人、男孩子、龍、兩個獸與十四萬四千人，認識屬靈爭戰、敵對神的權勢，以及聖徒的忠心。 \\

八 &
七碗、巴比倫與羔羊得勝 &
8-10 &
15:1-19:21 &
研讀七碗、巴比倫的傾倒、羔羊婚筵與基督再來；認識神的公義審判，以及世俗權勢終必敗亡。 \\

九 &
千禧年、撒但與最後審判 &
11 &
20:1-15 &
研讀撒但被捆綁、千禧年、歌革與瑪各，以及白色大寶座；比較前千禧年、無千禧年與後千禧年觀點。 \\

十 &
新天新地：被恩典帶入榮耀 &
12 &
21:1-22:21 &
研讀新耶路撒冷、神與人同住、生命河與生命樹；以新創造的盼望和「主耶穌啊，我願你來」作結。 \\

\end{longtable}

}